% ScholarPress starter manuscript for the "House preprint" profile.
% Upload this file as the main .tex and references.bib as the bibliography. Figures go in as supporting files
% with the same relative paths used in \includegraphics (for example figures/overview.png).
\documentclass[notice={Preprint. This version has not been peer reviewed.}]{scholarpress-preprint}
\usepackage[numbers,sort&compress]{natbib}

\title{A Starter Manuscript for ScholarPress Preprints}
\author{First Author}
\affiliation{School of Computing, Example University}
\author{Second Author}
\affiliation{Department of Software Engineering, Another University}
% \preprintdate{November 19, 2026}  % optional; defaults to the build date

\begin{document}
\maketitle

\begin{abstract}
This starter shows the elements the house preprint style supports: a title and author block with affiliations,
an abstract, numbered sections, displayed equations, a figure, cross-references, and linked references.
Replace the text, keep the structure, and rebuild.
\end{abstract}

\section{Introduction}\label{sec:intro}
Preprints let readers find work before formal publication~\cite{ginsparg2011arxiv}. Section~\ref{sec:method}
describes the method and Equation~\eqref{eq:loss} gives the objective.

\section{Method}\label{sec:method}
We minimise the regularised loss
\begin{equation}\label{eq:loss}
  \mathcal{L}(\theta) = \frac{1}{n}\sum_{i=1}^{n} \ell\bigl(f_\theta(x_i), y_i\bigr) + \lambda \lVert \theta \rVert_2^2 ,
\end{equation}
following standard practice~\cite{goodfellow2016deep}. Figure~\ref{fig:placeholder} is a placeholder figure.

\begin{figure}[t]
  \centering
  \fbox{\rule{0pt}{1.2in}\hspace{3in}}
  \caption{Replace this box with \texttt{\textbackslash includegraphics\{figures/overview.png\}}.}
  \label{fig:placeholder}
\end{figure}

\section{Conclusion}
State the contribution and its limits plainly.

\bibliographystyle{plainnat}
\bibliography{references}
\end{document}
